\section{انتگرال‌گیری عددی به روش‌های ذوزنقه‌ای، سیمپسون و کوادراتور گوسی}
در این بخش، هدف محاسبه عددی مقدار انتگرال تابع مشخص‌شده در صورت پروژه یعنی $ f(x) = e^{x^2} $ در بازه کران‌دار $ (0,1) $ است. از آنجا که این تابع فاقد پادمشتق برحسب توابع مقدماتی است، مقدار انتگرال آن صرفاً از طریق روش‌های تقریبی عددی محاسبه می‌شود. سه روش ذوزنقه‌ای\footnote{\lr{Trapezoidal Rule}}، سیمپسون\footnote{\lr{Simpson's Rule}} و کوادراتور گوسی\footnote{\lr{Gaussian Quadrature}} در محیط پایتون\footnote{\lr{Python}} روی این تابع اعمال شده‌اند.

\subsection{مراحل و نحوه پیاده‌سازی الگوریتم‌ها}
مراحل محاسباتی و الگوریتمی این بخش به شرح زیر است:
\begin{enumerate}
    \item \textbf{تعریف تابع هدف:} ابتدا تابع ریاضی $ f(x) = e^{x^2} $ با بهره‌گیری از تابع نمایی آرایه نام‌پای\footnote{\lr{NumPy}} یعنی \lstinline|np.exp| پیاده‌سازی می‌شود.
    \item \textbf{شبکه‌بندی بازه انتگرال‌گیری:} برای اعمال روش‌های مبتنی بر نیوتن-کوتس\footnote{\lr{Newton-Cotes formulas}} (ذوزنقه‌ای و سیمپسون)، بازه کل $ [0, 1] $ به کمک تابع \lstinline|np.linspace| به تعداد $ 100 $ زیربازه هم‌اندازه با مجموع $ 101 $ نقطه گره‌ای تقسیم‌بندی می‌شود.
    \item \textbf{اعمال قاعده‌ ذوزنقه‌ای و سیمپسون:} با استفاده از توابع اختصاصی و بهینه \lstinline|trapezoid| و \lstinline|simpson| از زیربسته محاسبات انتگرال سای‌پای\footnote{\lr{SciPy}} (\lstinline|scipy.integrate|)، ارزیابی عددی با دریافت فواصل نقشه گره‌ای انجام می‌پذیرد.
    \item \textbf{اعمال قاعده‌ کوادراتور گوس-کرونرود:} برای محاسبه انتگرال گوسی با بالاترین درجه دقت چندجمله‌ای، از تابع قدرتمند \lstinline|quad| استفاده شده است که علاوه بر مقدار انتگرال، تخمین دقیقی از خطای مطلق سقف محاسبات را نیز در خروجی برمی‌گرداند.
\end{enumerate}

\subsection{کد پایتون پیاده‌سازی شده}
کد کامل پایتون توسعه‌یافته برای محاسبات انتگرال‌گیری عددی به شرح زیر است:

\begin{latin}
\begin{lstlisting}
import numpy as np
from scipy.integrate import trapezoid, simpson, quad

# 1. Define the target function f(x) = e^(x^2)
def f(x):
    return np.exp(x**2)

# Define the integration intervals (from 0 to 1)
a = 0
b = 1

# Number of points for Trapezoidal and Simpson's methods
num_points = 101
x_intervals = np.linspace(a, b, num_points)
y_values = f(x_intervals)

print("--- Part 3 (a): Trapezoidal and Simpson's Methods ---")

# 2. Calculate using the Trapezoidal Rule
integral_trapezoidal = trapezoid(y_values, x_intervals)
print(f"Trapezoidal Rule Result: {integral_trapezoidal:.8f}")

# 3. Calculate using Simpson's Rule
integral_simpson = simpson(y_values, x_intervals)
print(f"Simpson's Rule Result:    {integral_simpson:.8f}\n")

print("--- Part 3 (b): Gaussian Quadrature Method ---")

# 4. Calculate using Gaussian Quadrature
integral_gaussian, abs_error = quad(f, a, b)
print(f"Gaussian Quadrature Result: {integral_gaussian:.8f}")
print(f"Estimated Absolute Error:    {abs_error:.2e}\n")

print("--- Summary & Comparison ---")
print(f"Difference (Gaussian vs Simpson):     {abs(integral_gaussian - integral_simpson):.2e}")
print(f"Difference (Gaussian vs Trapezoidal): {abs(integral_gaussian - integral_trapezoidal):.2e}")
\end{lstlisting}
\end{latin}

\subsection{خروجی دقیق برنامه در محیط ترمینال}
مقادیر عددی خروجی چاپ‌شده در ترمینال پس از اجرای اسکریپت فوق به صورت زیر ثبت شده است:

\begin{latin}
\begin{verbatim}
--- Part 3 (a): Trapezoidal and Simpson's Methods ---
Trapezoidal Rule Result: 1.46269705
Simpson's Rule Result:    1.46265175

--- Part 3 (b): Gaussian Quadrature Method ---
Gaussian Quadrature Result: 1.46265175
Estimated Absolute Error:    1.62e-14

--- Summary & Comparison ---
Difference (Gaussian vs Simpson):     3.02e-09
Difference (Gaussian vs Trapezoidal): 4.53e-05
\end{verbatim}
\end{latin}

\subsection{تحلیل و مقایسه جامع نتایج عددی روش‌ها}
با توجه به نتایج حاصل از خروجی سیستم، مقدار تقریبی انتگرال برای هر سه روش با دقت اعشاری بالا استخراج گردید. فرمول ریاضی هر یک از روش‌ها و تحلیل خطای آن‌ها به شرح زیر تبیین می‌شود.

\subsubsection{روش ذوزنقه‌ای}
در این روش، تابع در هر زیربازه با یک خط راست تقریب زده می‌شود. فرمول محاسباتی آن به صورت زیر است:
\[
I_T = \frac{h}{2} \left[ f(x_0) + 2\sum_{i=1}^{n-1} f(x_i) + f(x_n) \right].
\]
مقدار حاصل از این روش برابر با $ 1.46269705 $ به‌دست آمده است. این روش به دلیل استفاده از تقریب‌های مرتبه اول خطی، دارای خطای قطع کلی از مرتبه $ \mathcal{O}(h^2) $ است. تفاوت این روش با روش گوسی برابر با $ 4.53 \times 10^{-5} $ است که بیشترین میزان انحراف را در میان روش‌های مورد آزمایش نشان می‌دهد.

\subsubsection{روش سیمپسون}
در روش سیمپسون، تقریب توابع در هر دو زیربازه مجاور با استفاده از سهمی‌های درجه دوم انجام می‌گیرد. فرمول ریاضی این قاعده به صورت زیر است:
\[
I_S = \frac{h}{3} \left[ f(x_0) + 4\sum_{i=1,3,\dots}^{n-1} f(x_i) + 2\sum_{j=2,4,\dots}^{n-2} f(x_j) + f(x_n) \right].
\]
مقدار عددی حاصل از محاسبات سیمپسون برابر با $ 1.46265175 $ است. خطای قطع این روش از مرتبه $ \mathcal{O}(h^4) $ است. به دلیل انحنای سهموی توابع برازش شده، این روش انطباق بسیار بالایی با ساختار هموار تابع $ e^{x^2} $ دارد، به طوری که تفاوت خروجی آن با روش گوسی صرفاً برابر با $ 3.02 \times 10^{-9} $ است که نشان‌دهنده دقت فوق‌العاده بالای آن در تعداد بازه داده‌شده است.

\subsubsection{روش کوادراتور گوسی}
روش کوادراتور گوسی برخلاف دو روش قبل که از فواصل یکنواخت و ثابت استفاده می‌کردند، نقاط نمونه‌برداری $ x_i $ و وزن‌های متناظر $ w_i $ را به صورت بهینه و بر اساس ریشه‌های چندجمله‌ای‌های لژاندر\footnote{\lr{Legendre polynomials}} انتخاب می‌کند تا چندجمله‌ای‌های درجه بالا را دقیقاً انتگرال‌گیری کند:
\[
I_G = \sum_{i=1}^{m} w_i f(x_i).
\]
مقدار حاصل از این روش دقیق برابر با $ 1.46265175 $ با مرتبه خطای فوق‌العاده ناچیز و حد آستانه خطای مطلق محاسباتی معادل $ 1.62 \times 10^{-14} $ به‌دست آمده است. 

\subsubsection{نتیجه‌گیری مقایسه‌ای}
با در نظر گرفتن روش کوادراتور گوسی به عنوان دقیق‌ترین معیار انتگرال‌گیری عددی (به دلیل ساختار متغیر و بهینه نقاط ارزیابی)، مقایسه شاخص‌ها اثبات می‌کند که روش سیمپسون با داشتن خطای ناچیز مرتبه نهم، به طرز چشمگیری نسبت به روش ذوزنقه‌ای (با خطای مرتبه پنجم) ارجحیت دارد. دلیل این امر، نرم بودن و مشتق‌پذیری بی‌نهایت تابع تحت انتگرال است که باعث می‌شود فرمول‌های مرتبه بالاتر مانند سیمپسون و کوادراتور گوسی با سرعت بسیار بالاتری به سمت مقدار واقعی انتگرال همگرا شوند.
